\section[第10章\q 现在就过富足的生活]{第10章\\ 现在就过富足的生活}
\label{sec:ch10}

周五早上并不顺遂。因为打定主意要带午餐，柔伊洗手下厨，她从自己编辑的杂志中选了一款简单的地中海料理食谱。结果她没有成功，反而煮出一堆难以下咽的烧焦青菜，让她的心情更糟糕。

前一晚她还做了一个恶梦，梦中恐怖的跑步机又出现了，这回跑步机悬在一个滚烫的岩浆上方。她必须加快脚步，不然就会掉下去，把自己烤焦。她可以感受到下发出吱吱声响，不停冒出来的热气。热灰落在她的头发上、脸上。最后，终于在自己的尖叫声中惊醒了。

躺在凌晨3:00的黑暗中，她做了一个决定。

如果等一下，洁西卡介绍的经纪公司打电话来，她就说好。

处理好春季特刊内容后，她会邀芭芭拉共进午餐，不是在公司的餐厅，而是去外头某家餐厅。

她也不期待这样的餐叙，只是这是必须做的事。芭芭拉可能会说，她想要的高薪不是财务问题唯一且万能的答案。但是，拜托现实一点，收入大幅提升也没什么不好。

柔伊收拾好包包，走出客厅，发现外头在下大雨，又折回来取了雨伞，才往L地铁路走去，一路上小心避开地上的积水。这样出门时间就比平常晚了一点，她想今天就不要去咖啡厅了，那些百万退休金的神奇事她也听得够多了。她应该没空去享用拿铁和马芬了，直接转车去公司。

可是在走在街上时，她的编辑小脑袋却又忍不住开始回想亨利说过的话。

在做杂志时，她也常常遇到同样的情形，许多好点子挤在一起，令人眼花撩乱，无法找到重点。有时作者在一篇文章中，有太多好点子；有时他们写着写着，偏离了原本的重点；有时候，一个很好的想法，却没有发展完整，无法得出一个关键结论。

亨利所说的话，关键的结论是什么呢？

快到地铁站口时，柔伊突然停住了脚步。她对自己突如其来停住，因而挡到其他乘客行径连声道歉。

亨利怎么说的？要富有，要能财务自由，其实很简单，只要做到三步骤。我称之为财富自由的三个秘密步骤。

先付钱给自己、养成习惯……那才两个啊！

第三个秘密，是什么？

柔伊直接掉头，往海伦娜咖啡厅走过去。

\subsection{马上为自己创建梦想帐户}

10分钟后，柔伊的雨伞就卷好放在咖啡厅门口的雨伞架子里，而她自己则坐在亨利对面的椅子上。

“喔，”亨利说。“你要问第三个秘密啊，没问题。”他往后靠了靠，手放在膝盖上。“让我们先来聊一聊，什么才是最重要的。”

“好的。”柔伊说。“来聊吧！快点告诉我，什么才是最重要的。”

“不，”亨利笑着摇头。“不是这样的，是你告诉我。”

“我不太清楚你的意思。”

“我们之前谈过的退休帐户，好像那是一件很重要的事……但是究竟有多重要？我的意思是，对你来说，退休金到底有多重要。是的，你知道这笔钱拿到是约半世纪之后了。不像于我，一个七十多岁的老人，我正靠这个过生活呢！而对你这样的年轻人来说，这是不是太不真实？我说得对吗？”

这点他其实说得没错。就像芭芭拉提过的，那些在老年过得贫乏的女人，光是想像就让柔伊吓坏了。40年后简直好像要等上一辈子才会到似的。

“好吧。那我们现在先不管退休之后。我们先说说你现在的生活，从今开始的四、五十年生活。你都做些什么梦呢？”

“我的梦？”柔伊耸了耸肩膀，并心想：喔，我想你不会想知道的。

“不是说你做的恶梦。”亨利温和的说，好像他读懂了她的念头一样，“我是说你的梦想。说说你真正想做的事。”

“我想学摄影，拍出精采的照片。”柔伊想都没想，便脱口而出。

“摄影课程。”亨利点头，“不错。”

“也不是什么拯救世界的大抱负。”柔伊说。

亨利若有所思的抬起头来说：“不用觉得不好意思。”他说，“大的梦想未必好。梦想就是想做的事。有时候，最简单的反而最强大，也更容易实现。比方说，上摄影课。你想上摄影课，那为什么不直接去上？”

柔伊才要开口说话，亨利竖起食指说：“等一下，有个条件，那就是不可以说‘我付不起学费’。”

“好的。”柔伊说。她想了一下后说：“因为学费太贵。”

亨利忍不住笑了。

事实上，她已经想去上这个当地的课程好几年了。也不是很贵，不超过600美元。可是她竟然凑不出600元现金。

“好吧。”亨利说：“我们来看看这点。你的薪水有没有自动转到401帐户的设置？”看到柔伊犹豫的表情，亨利说，“那就是没有了。”

“是的，我没有设置。”柔伊承认，“但是我正在考虑。”

亨利低下头，并对柔伊露出很严肃的表情：“好，我不会评论这件事。”

柔伊脸上的笑容很无辜。

“我们先把退休帐户放在一边，\textbf{我觉得你需要一个梦想帐户}。和退休帐户完全分开，梦想帐户就是上学基金，去学想学的东西。那么，我们且称之为柔伊的摄影课帐户吧，有钱转帐进去，比方说，1个月转个100美元，1天只要3.5美元。”（也就是再一杯拿铁的钱，柔伊想。）“你刚刚说学费是多少？”

“600美元上下。”柔伊说。

“那好。就1个月100美元，6个月后去上摄影课，梦想达成。再来，你还想做什么？”

柔伊皱眉。她也不知道为什么，有时候她的脑袋就是打了结。“我……”她看着亨利说，“就是脑中一片空白。”

“试试这一招。”他说，“闭上眼睛，深呼吸，吸气、吐气。”

柔伊深吸了一口气，再吐出来。

“很好喔，”亨利说，“现在再想一想，生活中有什么你最想做的。”

柔伊再次深吸一口气，缓缓吐出来。

她脑中浮现出7岁的某一天，她坐在车子后座，父母在前座。他们一路往北，要去缅因州的海岸。

他们到达之后，在海边散步。低矮的灌木，飞翔的海鸥，海岸巨大的岩石和冰冷的海水。他们从海边回到旅馆用餐，好美味的蓝莓。还有派，当然也是蓝莓口味的，那是她一生中吃过的最好的蓝莓派。

“哇，”她感叹着，闭上眼睛。有好些年，她都没有想起这趟旅程。她喃喃对亨利说起了儿时的回忆。

她记得他们三个搭捕虾船出海。在墨绿色的水波上前行。船长让她掌舵时，木质舵柄的握感。她一五一十都对亨利说了。

\subsection{人生最重要的事：实现梦想}

“后来，我再没有搭过船。”柔伊说。

“那么告诉我，柔伊，”她听见亨利温和的说。“你当时在缅因州海面的小船上，是什么感觉？这趟旅程中哪些部分是你喜欢的？”

柔伊睁开眼睛，碰上了亨利的眼神，他的眼中充满期待。

“感觉好像在冒险。”她说。“去哪里都好。好像在天空飞翔一样，非常自由。”她停顿了一下，又重复说：“非常自由。”

她闭上眼睛，回想着这句话。

自由塔\footnote{The Freedom Tower：世界贸易中心一号大楼，位于美国纽约曼哈顿下城的摩天大楼，为世界贸易中心在911事件遭损毁后重建的建筑之一。原称自由塔，于2009年更名，但仍有许多纽约人称呼其旧名。}。

她每天从办公室的小餐厅往外看，看到的景象是自由女神。

“我想，可能那就是我想要的。”她小声嘀咕，“不单是摄影课，还有自由的感觉。我只是想知道，\textbf{我可以做我想做的事，去我想去的地方，而且是在我想去的时候，说走就走}。”柔伊睁开眼睛，脸涨红了，“我知道这听起来很自我，很不现实。”

亨利眼睛都没眨一下，说：“我不同意你这么说。对我来说，这很合理。我们每个人来地球上，都有一个特别的使命。你想拥有实现自己使命的自由，对我来说，很正常的。”

柔伊微微点头。“也对。我明白你说的。”

“现在告诉我，”亨利说，“当你搭着小船在缅因州海面，感受到自由的时候，带给你什么感受？”

柔伊再次闭上了眼睛，回忆再次来到船上，一个词从她口中跳出来：“冒险旅程。”她睁开眼，看着亨利。“我过去从来没想到，但这就是我想要的。冒险的自由。去看过没看过的景致，去没去过的地方。”

亨利点头后说：“你在一个什么样的地方工作？”

柔伊疑惑的看了亨利一眼，就笑了起来。“这个问题问得好。”身为旅游杂志的编辑，她的工作就是修润别人出游的文本。

别人的冒险。

“我可以再问得详细一点吗？”亨利说，“是哪一类的冒险？可以带给你什么感受的冒险？”

柔伊再次闭上眼睛。如果她可以去任何想去的地方，做任何想做的事，那该是哪里？“不是跳伞，也不是摩托车越野，”她说，“也不是爬山。就是想去一个地方，看世上最漂亮的风景。”

她停顿了一下，继续说：“用想的就很激动，对不对？想到即将看到的美景。”

说完，她睁开了眼睛。

亨利拿出笔，翻开笔记本新的一页。写下三个字：

{\kaishu 自由、冒险、美景}

“柔伊，你知道，为什么许多人没办法存到钱吗？或者，他们即便存了，也存了很少，派不上大用处？就是因为他们没有看到这一点。”

“这些，”亨利点头示意写在纸上的字，“这就是关键。人们常常说，更大的房子、度假小屋，或者高薪。但这些其实并不重要，这几点才是最重要的。

“\textbf{因为这是你的梦想}。不管是短期的，比方一堂摄影课，还是长期的，如环游世界，\textbf{这些需要去实现的梦想，才是人生中最重要的事}。他们比生命更重要，就像是氧气，缺少了梦想，人生就会窒息。”

“这个清单或许还不完整。毫无疑问，你会再做编辑、修改和增补。或许你可以说，这就是你人生的价值。对你来说，这是最最重要的。”

亨利朝着餐巾纸点点头，“现在的问题是，你每天做的事情和选择，让你离梦想更近了吗？在花钱时，也是把梦想放在第一位了吗？”

柔伊想了想，说道：“你想说的是，让我自己做了最想做的事吗？”

亨利说：“没错，那正是我想说的。”

柔伊看着纸上的几个大字，她不晓得为什么又想起了老妈说的话：要为已经拥有的而开心。她有做到吗？

她看了看亨利。“你介意我请教一个问题？”

亨利笑着说：“请问。”

“什么事情能带给你最单纯、最雀跃的喜悦感？”

亨利往后靠了靠，凝视柔伊好一段时间，然后微微点头。

“这是个好问题。”他站起身来说，“我们边走边聊？”他们缓慢起身，亨利同时开始侃侃而谈。

\subsection{真正的富足，不在金钱，而在生命}

“36年前，有位好友问过我同样的问题。我当时很惊讶，因为我从来没有问过自己这个问题，所以一时竟然答不出来。

“我当时并不快乐。我很年轻，在一家知名的设计公司上班，有钱途也有前途。我喜欢我的工作，也喜欢一起工作的人。但这就是我真正想要的吗？”他缓缓摇头，“我必须承认，这些都不是。\textbf{我每天花许多时间在工作上，但是却没有让我离想要的生活更近}。

“于是我找老板谈了谈，请了一个长假。打包了行李，订了机票，就去了欧洲。一开始，我对老板请的长假是几周的时间。我也称之为‘高调的休假研习’。”他呵呵笑了，“可是后来的结果是，我再也没有回去。

“我朋友都说我疯了，好好一份工作，被我搞丢了。我也知道，他们说得没错。有好几年的时间，我都不停地告诉自己说，‘有一天，亨利，你会环游世界，看到这个星球上最漂亮最迷人的所在，亲眼看到他们。’

“所以，我就去做了。那种工作到退休然后再去环游世界的模式，对我来说，突然失去了意义。我回到家之后，就去申请了小额贷款，在喜欢的社区租了一处小地方。”

“开始卖咖啡。”柔伊说。

亨利点头，“于是开始了先付钱给自己的人生，拥有了属于自己的人生。也没过几年，我就买下了这栋大楼。从第一趟欧洲旅行之后，每年我都用6周的时间去看世界。过去的36年里，我去过了一百多个国家。”

柔伊突然恍然大悟，她心底一颤。

她低声说：“那些照片，都是你的杰作。”

亨利看着柔伊，笑道：“就像我对你说过的，你喜欢的那张也是我自己最喜欢的。”

他们在黎明的米克诺斯岛照片前方停下了脚步，两个人一起看着画面。

“拍下这张照片时的感受，直到今天仍历历在目。”亨利的声音变得很遥远，且很不平静。

“按下快门后，我就转身放下相机。然后单膝跪下求婚。”

“她说了‘我愿意’。”柔伊低声问。

“是的，她说了‘我愿意’。”亨利说。

“是海伦娜吗？”

亨利笑了：“海伦娜。就像特洛伊的海伦一样，全希腊最美丽的女人。拍照地点也是我们初识的地方，我们相遇在几周之前。这就是我的第一次旅行，我高调的学术休假。她和我一起回到了美国……”他停顿下来，朝着照片墙伸展开手臂，好像在说：她现在就在这里。

在海伦娜咖啡厅里。

“我生命的挚爱。”他说，“此后，对她，我也只说，是的，好的，我愿意。”

柔伊终于明白为什么特别受到这张照片的吸引了。不单单因为景美，更有浓浓的人情，爱和无尽的可能，在瞬间爆发出来，如画面上的金色光芒。

不单这一张照片，挂在墙壁上的所有照片，都是亨利生命中的重要时刻，永恒被留存的时刻。

柔伊随后又领悟出，她一直觉得亨利身上有种特质吸引着她，这种特质也吸引着拜伦、乔琪亚还有芭芭拉，好像磁力一般，或许就是所谓的魅力吧，但又不止于此。那是一种平静的满足。现在他生命的重要瞬间都展现在她的眼前，栩栩如生。

不是他老派的绅士风度，不是他带着异国情调，甚至也不是他的智能。

而是他的富足。

\textbf{不是金钱上的富有，而是生命的富足。}

“所以，这就是第三个秘密吗？”她问。

亨利笑着。“是的。没有第三个秘密，前两个就没有实现的动力了，因为你会不想照着做。”

他把笔记本往前翻，翻到写着两个秘密的那一页，加上了第三个：

{\kaishu 3.现在就过富足的生活}

“如果说前两个秘密“先付钱给自己”和“养成习惯”是具体作为，那么第三个就是动机。想一想，什么对你来说最重要，然后甘心为此调整自己的生活。

“现在就过富足生活。而不是等到遥远的未来。今天，现在。”
